%% ---------------------------------------------------------------------------------------
\section{Interpretations and conclusions}
\label{sec:conc}

\begin{enumerate}[leftmargin=*]
\item \textbf{The tree is a constructive proof of interpolation, and a microscope on it.} The note's
  construction carries over to neural nodes almost unchanged: hinge-loss nodes, closed-form
  connection weights, and an isolation node to guarantee termination. Because every node knows which
  samples it is responsible for, the tree makes visible \emph{who} memorises \emph{what}: the root
  fits the regular structure, first-level children the atypical members of each class, deeper nodes
  individual cluttered or mislabelled-looking photographs, and -- with label noise -- most of the first-level children's positives are corrupted labels.
\item \textbf{The price of 100\% training accuracy depends on how locally one memorises.}
  Interpolation is not harmful or benign in itself. Learned correctors generalise the memorised
  points to whole regions -- deep ones literally fire on ``anything that is not a known negative'' --
  and cost up to 5 test points; isolation nodes (max-margin caps around single
  points) are a lookup table that never fires on test images and costs nothing, at the price of one
  hidden unit per memorised image. This is the tree-level version of the benign-overfitting picture
  \cite{bartlett,feldman}: memorising the long tail is harmless when it is spiky and local.
\item \textbf{Double descent is a property of the root; the tree jumps over the classical regime.}
  The peak sits exactly where the root first interpolates, with and without noise. Growing children
  reaches interpolation with far fewer parameters per node than widening, but it reaches the same
  over-parameterised error level, never the classical optimum.
\item \textbf{Features decide everything that matters.} Pixels give large trees that generalise poorly
  (60\% on Cats vs Dogs); a trunk trained on the same data hides memorisation inside the trunk;
  a set of trunks with per-node selection gives the smallest and best trees, because a node's errors
  are best separated in a feature space other than the one that produced them.
\item \textbf{Small and interpolating are at odds with good test accuracy.} The smallest
  interpolating network is the MLP at the threshold width -- the double-descent peak. Within a fixed
  tree, decision-preserving pruning removes 20--32\% of the weights for free, and cost-complexity
  pruning shows a near-linear exchange of 0.3--0.5 test points per training point.
\end{enumerate}

\paragraph{Challenges encountered.} (i) Child problems are extremely imbalanced and ``inside-out'';
without class weights the child learns ``never fire'', and linear children almost never make progress.
(ii) The connection weights of the LP are badly scaled (hundreds) -- the closed form is both exact and
better conditioned. (iii) A trunk trained on the same images inflates the root's training accuracy.
(iv) On a 2-core CPU, the trunks have to be frozen and features cached; oversubscribing the cores
slowed PyTorch down five-fold.

\section{Future directions}
\begin{itemize}[leftmargin=*]
\item \textbf{Local neural correctors.} Gate each learned corrector by an isolation cap around its
  own positives (fire only if both agree): it would keep the compactness of the network and the
  test-time locality of the lookup table.
\item \textbf{Multi-head detectors.} Replace the $K$ detectors $D_k$ by one network with $K$
  outputs and a shared hidden layer, trained jointly with the Table~II targets of all classes.
\item \textbf{Fine-tuning the trunk with the tree.} Back-propagate the corrector losses into the
  last trunk layers, so that the features themselves are reshaped to make the errors separable.
\item \textbf{Margin-based growth past interpolation.} Keep growing children for samples whose margin
  is below $\gamma>0$; by analogy with boosting~\cite{schapire} this may continue to improve test error
  after the training error is zero.
\item \textbf{Minimal-complexity nodes.} Train every node with the minimal complexity machine
  objective~\cite{mcm} or with the zero-norm~\cite{weston}, as the note suggests, to obtain the sparsest
  node that solves its sub-problem.
\end{itemize}
